\documentclass[]{article}

\usepackage{graphicx}
\usepackage{booktabs, tabularx, wrapfig}
\usepackage[hidelinks]{hyperref}
\hypersetup{hidelinks}
\usepackage{xcolor}
\usepackage[normalem]{ulem}
\usepackage{multirow}
\usepackage{longtable}
\usepackage{amsmath, amssymb}
\usepackage{enumitem, array}

\begin{document}

\textbf{FACULDADE DE TECNOLOGIA DE SÃO PAULO -- FATEC-SP}

\textbf{EDUARDO DE MATTOS BORTOTO}

\textbf{Desenvolvimento de um Laboratório de Detecção de Intrusão:
Comparando Machine Learning e Deep Learning na Identificação de Ameaças
em Redes}

São Paulo

2026

\textbf{FACULDADE DE TECNOLOGIA DE SÃO PAULO -- FATEC-SP}

\textbf{EDUARDO DE MATTOS BORTOTO}

\textbf{Desenvolvimento de um Laboratório de Detecção de Intrusão:
Comparando Machine Learning e Deep Learning na Identificação de Ameaças
em Redes}

Trabalho de conclusão do curso apresentado à FATEC-SP, como exigência
parcial para a obtenção do grau de Tecnólogo em Análise e
Desenvolvimento de Sistemas, sob orientação da Professora Mestre Grace
Anne Pontes Borges

São Paulo

2026

SUMÁRIO

Contexto Atual

A segurança da informação atravessa, na atualidade, um momento de
ruptura paradigmática. Durante décadas, a proteção das redes
corporativas fundamentou-se em uma lógica reativa, onde a defesa
dependia quase exclusivamente do conhecimento prévio das ameaças.
Ferramentas clássicas, como firewalls e sistemas de detecção de intrusão
(IDS) baseados em assinaturas, operavam de maneira similar a um sistema
imunológico que só reconhece patógenos já catalogados. Se o tráfego de
rede contivesse uma sequência de bytes conhecida como maliciosa, o
bloqueio era imediato; caso contrário, o acesso era permitido. No
entanto, a transformação digital acelerada, impulsionada pela adoção
massiva da computação em nuvem e pela expansão da Internet das Coisas
(IoT), expôs a fragilidade desse modelo. O perímetro de segurança, antes
bem definido pelas paredes do data center, dissolveu-se, e o volume de
dados trafegados atingiu escalas que tornam a análise manual humana
impraticável.

O cenário agravou-se com a sofisticação das táticas ofensivas.
Cibercriminosos utilizam hoje ataques de "Dia Zero" (Zero-Day), que
exploram vulnerabilidades inéditas para as quais não existem assinaturas
de defesa, tornando invisíveis as invasões aos olhos dos sistemas
tradicionais. Somado a isso, a onipresença da criptografia em protocolos
como HTTPS e TLS 1.3, embora essencial para a privacidade, criou um
"ponto cego" para a defesa, impedindo a inspeção profunda do conteúdo
dos pacotes. Nesse contexto, a Inteligência Artificial (IA) emerge não
como um luxo, mas como uma necessidade operacional. A transição para
sistemas baseados em anomalia comportamental permite que a defesa
aprenda o padrão "normal" da rede e identifique desvios sutis,
possibilitando a detecção de ataques nunca antes vistos,
independentemente de sua assinatura ou da criptografia utilizada.

Para profissionais que atuam na linha de frente da segurança ofensiva,
como é o caso dos pentesters, compreender essa evolução na defesa é
vital. A dicotomia tradicional entre quem ataca (Red Team) e quem
defende (Blue Team) está dando lugar a uma abordagem integrada,
conhecida como Purple Team. Entender a engenharia por trás dos
algoritmos de detecção modernos não serve apenas para construir defesas
melhores, mas também para aprimorar as técnicas de auditoria, garantindo
que os testes de segurança reflitam a realidade de um ambiente
monitorado por inteligência artificial.

\begin{enumerate}
\def\labelenumi{\arabic{enumi}.}
\item
  Objetivo
\end{enumerate}

\begin{quote}
O objetivo central deste trabalho é desmistificar a aplicação de
Inteligência Artificial na segurança defensiva através da construção de
um protótipo funcional de um Sistema de Detecção de Intrusão (NIDS). Em
vez de focar apenas na teoria abstrata, o projeto propõe o
desenvolvimento de um laboratório prático onde será possível treinar,
testar e comparar a eficácia de diferentes gerações de algoritmos na
distinção entre tráfego legítimo e malicioso.

De maneira mais específica, o estudo busca avaliar se a complexidade
computacional das modernas redes neurais de Deep Learning (como
Autoencoders) realmente se traduz em um ganho de segurança significativo
quando comparada a algoritmos de Machine Learning mais tradicionais e
leves, como o Random Forest. O trabalho envolverá todo o ciclo de vida
de um projeto de ciência de dados aplicado à segurança: desde a coleta e
limpeza de bases de dados reais de tráfego de rede, passando pela
engenharia de atributos para selecionar as características mais
relevantes do tráfego, até a implementação e validação dos modelos.
\end{quote}

\begin{enumerate}
\def\labelenumi{\arabic{enumi}.}
\setcounter{enumi}{1}
\item
  Justificativa
\end{enumerate}

\begin{quote}
A motivação para este trabalho emana da intersecção entre o interesse
acadêmico e a minha trajetória profissional na Caixa Econômica Federal.
Atuando como Pentester, minha rotina envolve a emulação de táticas
adversárias para identificar vulnerabilidades preventivamente. Contudo,
a consolidação de ecossistemas críticos, como o Pix e o Open Finance,
ampliou significativamente a superfície de ataque e a velocidade das
transações financeiras. Nesse cenário, as ferramentas de defesa manuais
tornam-se obsoletas frente à sofisticação das fraudes automatizadas,
exigindo uma postura de segurança cibernética mais autônoma, preditiva e
inteligente.

Desenvolver este laboratório é um passo estratégico para minha carreira
em direção ao perfil de Purple Team. Ao construir meu próprio motor de
detecção baseado em IA, deixo de tratar as ferramentas de defesa como
"caixas pretas". Isso me permitirá compreender exatamente quais
comportamentos em um teste de intrusão disparam os alertas dos
algoritmos, tornando meus relatórios de pentest muito mais valiosos para
a equipe de defesa do banco. Além disso, o domínio de Python e
bibliotecas de Data Science é hoje um requisito fundamental para a elite
da cibersegurança, permitindo a criação de automações que vão muito além
dos scripts prontos.

Do ponto de vista social e institucional, contribuir para o estudo de
defesas mais eficientes em bancos públicos transcende a questão técnica.
Trata-se de colaborar para a integridade do sistema financeiro nacional
e para a proteção do patrimônio de milhões de cidadãos brasileiros que
dependem da disponibilidade e confidencialidade dos serviços digitais da
Caixa.
\end{quote}

\begin{enumerate}
\def\labelenumi{\arabic{enumi}.}
\setcounter{enumi}{2}
\item
  Hipóteses
\end{enumerate}

\begin{quote}
A condução do experimento prático será norteada por duas hipóteses
principais, formuladas com base na literatura técnica e na experiência
empírica de campo.

A primeira hipótese sugere que os modelos de Deep Learning não
supervisionados, especificamente os Autoencoders, apresentarão uma
capacidade superior de generalização frente a ataques desconhecidos
(Zero-Day) em comparação aos modelos supervisionados clássicos.
Acredita-se que, por aprenderem a reconstruir o padrão do tráfego normal
em vez de memorizarem assinaturas de ataques específicos, essas redes
neurais conseguirão sinalizar anomalias inéditas com maior precisão,
superando a limitação do Random Forest, que tende a ter bom desempenho
apenas contra padrões que já "viu" durante o treinamento.

A segunda hipótese aborda a eficiência operacional, postulando que é
possível construir um sistema de detecção robusto utilizando apenas os
metadados do fluxo de rede (como tempo entre pacotes, tamanho do
cabeçalho e contagem de flags), sem a necessidade de descriptografar o
conteúdo da comunicação. Se confirmada, essa hipótese valida a
viabilidade de implementar monitoramento inteligente que respeita a
privacidade do usuário e a criptografia ponta a ponta, mantendo a
performance necessária para redes de alta velocidade sem introduzir
latência excessiva.
\end{quote}

\begin{enumerate}
\def\labelenumi{\arabic{enumi}.}
\setcounter{enumi}{3}
\item
  Metodologia da Pesquisa
\end{enumerate}

\begin{quote}
O trabalho adotará uma metodologia de natureza aplicada e experimental,
estruturada em duas fases distintas mas complementares: uma
fundamentação teórica focada e um estudo de caso prático em laboratório.

Na fase inicial, será realizada uma revisão bibliográfica seletiva,
priorizando artigos técnicos e publicações recentes (2020-2025) que
abordem a aplicação de Machine Learning em cibersegurança. O objetivo
não é exaurir a literatura acadêmica, mas sim mapear o estado da arte
das técnicas de detecção de anomalias e entender como grandes players do
mercado estão substituindo as regras estáticas por modelos
probabilísticos.

A fase principal consistirá na montagem do laboratório experimental em
um ambiente controlado. Utilizarei a linguagem Python, amplamente
adotada tanto na comunidade de segurança ofensiva quanto na de ciência
de dados, juntamente com bibliotecas especializadas como Pandas para
manipulação de dados, Scikit-learn para os modelos clássicos e
TensorFlow/Keras para as redes neurais. O "combustível" para os testes
será o dataset CIC-IDS2017, escolhido por sua fidelidade em representar
uma rede moderna com tráfego real de ataques como Brute Force, DDoS e
Infiltration, superando bases obsoletas como o KDD-99.

O procedimento experimental seguirá um fluxo lógico: primeiramente, os
dados brutos passarão por um processo de limpeza e normalização. Em
seguida, os modelos serão treinados e submetidos a testes cegos com
tráfego misto. A avaliação final não se limitará a uma simples taxa de
acerto, mas focará na análise da Matriz de Confusão, permitindo mensurar
o equilíbrio entre a capacidade de detectar intrusões reais e a
habilidade de evitar os custosos falsos positivos, simulando assim os
desafios reais de um Centro de Operações de Segurança (SOC).
\end{quote}

\begin{enumerate}
\def\labelenumi{\arabic{enumi}.}
\setcounter{enumi}{4}
\item
  Referências Bibliográficas
\end{enumerate}

\begin{quote}
SHARAFALDIN, I. et al. Toward Generating a New Intrusion Detection
Dataset and Intrusion Traffic Characterization. \emph{ICISSP}, 2018. (O
artigo oficial do dataset que será usado no laboratório).

ALI, M. L. et al. Deep Learning vs. Machine Learning for Intrusion
Detection in Computer Networks: A Comparative Study. \emph{Applied
Sciences}, v. 15, 2025. (Base para a comparação entre os algoritmos).

LIAO, H. et al. A Survey of Deep Learning Technologies for Intrusion
Detection in Internet of Things. \emph{IEEE Access}, 2024. (Contexto de
IoT e modernização da defesa).

FEBRABAN. \emph{Pesquisa Febraban de Tecnologia Bancária 2024}. (Dados
sobre o investimento dos bancos em segurança e IA).

DENNING, D. E. An Intrusion-Detection Model. \emph{IEEE Transactions on
Software Engineering}, 1987. (O clássico que definiu o que é um IDS).

SOPHOS. \emph{The State of Ransomware 2024}. (Estatísticas atuais sobre
ataques que justificam a necessidade de melhor defesa).
\end{quote}

\end{document}
